%\usepackage{mathtools}
%\usepackage{microtype}
%\usepackage{titlesec}

%\titleformat{\subsubsection}{\normalfont\bfseries\small}{}{0em}{}
%\titlespacing{\subsubsection}{0pt}{3pt}{1pt}


\documentclass[landscape, 8pt]{extarticle}

\usepackage{../../preamble}
\usepackage[separate]{../../rss/thmboxes_v5}


\begin{document}
	
	
	\fontsize{5pt}{5pt} \selectfont
	
	%Reduce horizontal padding in 'tabular' environment, make rows taller for legibility
	\setlength\tabcolsep{4.5pt}
	\setlength\extrarowheight{2pt}
	
	
	\setlength{\abovedisplayskip}{0pt}
	\setlength{\belowdisplayskip}{0pt}
	\setlength{\abovedisplayshortskip}{0pt}
	\setlength{\belowdisplayshortskip}{0pt}
	
	
	\begin{multicols*}{4}
	\raggedcolumns
	%Original content by Louis and Leon's friend John's friend (dunno who), converted to LaTeX from PDF by Claude Sonnet 4.6 (many GPT versions failed and got weirdly sycophantic, Llama also failed, Claude did it first try), then compacted and condensed by Louis
	\vspace*{-14pt}
	\subsubsection*{CMM Exam Sheet 2025/26 - Non ii, non vidi, non vicī}
	\vspace{-6pt}
		
		\textbf{Thm:} (Stokes). Let $\Sigma$ be a smooth oriented surface in $\mathbb{R}^3$ with boundary $\partial\Sigma \equiv \Gamma$. If a vector field $F(x,y,z) = (F_x(x,y,z), F_Y(x,y,z), F_z(x,y,z))$ is defined and has cont. 1$^{st}$ order partial derivatives in a region containing $\Sigma$, then
		$
		\iint_\Sigma (\nabla \times F)\cdot d\Sigma = \oint_{\partial\sigma} F\cdot d\Gamma.
		$
		More explicitly, the equality says that
		$
		\iint_\Sigma\!\!\left(\!\left(\frac{\partial F_z}{\partial y}-\frac{\partial F_y}{\partial z}\right)dydz + \left(\frac{\partial F_x}{\partial z}-\frac{\partial F_z}{\partial x}\right)dzdx + \left(\frac{\partial F_y}{\partial x}-\frac{\partial F_x}{\partial y}\right)dxdy\right)
		= \oint_{\partial\Sigma}\!\!\left(F_xdx + F_ydy + F_zdz\right).
		$
		
		\textbf{Dfn:} (SVCDE Identities). $\bm{a} \times \bm{b} = \left\lVert \bm{a}\right\rVert \left\lVert \bm{b}\right\rVert \sin( \theta) \bm{n}$, $\bm{n}$ norm to $\bm{ab}$-plane
		\vspace{-10pt}
		\begin{multicols}{3}
			\begin{itemize}
				\item $\nabla\cdot(\nabla\times F)=0$
				\item $\nabla\times(\nabla F)=0$
				\item $\mathbf{x} \times \mathbf{x} = 0$
				\item $\mathbf{x} \times \mathbf{y} = - (\mathbf{y} \times \mathbf{x})$
				\item $\mathbf{x} \cdot \mathbf{y} \le \left\lVert \mathbf{x}\right\rVert \left\lVert \mathbf{y}\right\rVert$
				\item $\left\lVert \mathbf{x} + \mathbf{y}\right\rVert \le \left\lVert \mathbf{x}\right\rVert + \left\lVert \mathbf{y}\right\rVert$
				\item $\nabla(|\mathbf{x}|)=\frac{\mathbf{x}}{|\mathbf{x}|}$
				\item $\nabla\!(\frac{k}{|\mathbf{x}|})=-\frac{k\mathbf{x}}{|\mathbf{x}|^3}$
			\end{itemize}
		\end{multicols}
		\vspace{-10pt}
		\textbf{Thm:} For $C$ a curve $P \to Q$, $\int_C \nabla F\cdot ds = F(Q)-F(P)$
		
		\hrule
		\vspace*{-3pt}
		\subsubsection*{\tiny{1. Newtonian Mechanics}}
		\vspace*{-6pt}
		
		\textbf{Dfn:} The universe, to a good approx, looks like $\mathbb{R}\times\mathbb{R}^3$, where the 1$^{st}$ factor $\mathbb{R}$ represents \textbf{time}, and the 2$^{nd}$ factor $\mathbb{R}^3$ represents \textbf{space}.
		
		\textbf{Dfn:} A point $(t,\mathbf{a})\in\mathbb{R}\times\mathbb{R}^3$ is called an \textbf{event}. Two events $(t,\mathbf{a})$ and $(t',\mathbf{b})$ are \textbf{simultaneous} iff $t=t'$. The distance between simultaneous events is given by the norm $|\mathbf{a}-\mathbf{b}|$.
		
		\textbf{Dfn:} A particle's trajectory through the universe is given by the \textbf{worldline}. This is a graph
		$
		\mathbf{x}:\mathbb{R}\to\mathbb{R}^3, t\mapsto \mathbf{x}(t),\text{of form}\{(t,\mathbf{x}(t))\in\mathbb{R}\times\mathbb{R}^4\},
		$
		in the Newtian universe.
		
		\textbf{Dfn:} The \textbf{configuration space} of a many body system is given by the locations of each particle in space, so $\mathbb{R}^3\times\mathbb{R}^3\times\cdots\times\mathbb{R}^3 = \mathbb{R}^N$, $N\equiv 3n$. The \textbf{state space} of such systems is given by the collection of positions and velocities of each of the particles $(\mathbf{x},\dot{\mathbf{x}})$. Thus has dim $= 6n$.
		
		\textbf{Dfn:} For $\mathbf{x}_a(t):\mathbb{R}\to\mathbb{R}^3$ a worldline of particle $a$, the collection of these worldlines defines a curve
		$
		\mathbf{x}(t)=(\mathbf{x}_1(t),\mathbf{x}_2(t),\ldots,\mathbf{x}_n(t)),
		$
		in the configuration space.
		
		\textbf{Princ:} (Newt's Principle of Determinancy). The motion of any mechanical system is uniquely determined by its initial state, the set of all its constituent positions and velocities at a given instant in time. Equivalently the worldline $\mathbf{x}(t)$ in the configuration space is determined by $\mathbf{x}(0)$ and $\dot{\mathbf{x}}(0)$ for all $t\in\mathbb{R}$ or at least for all $t\in[0,\tau)$ for some $\tau>0$.
		
		\textbf{Dfn:} (Newt's Eqn). By Newt's PoD there must be a relationship of form
		$
		\ddot{\mathbf{x}}(t)=\Phi(\mathbf{x},\dot{\mathbf{x}},t)
		$
		for some $\Phi:\mathbb{R}^N\times\mathbb{R}^N\times\mathbb{R}\to\mathbb{R}^N$. Let $v:\mathbb{R}\to\mathbb{R}^N$ then the system
		$
		\begin{psmallmatrix*}\dot{\mathbf{x}}(t)\\ v(t)\end{psmallmatrix*} = \begin{psmallmatrix*}v(t)\\ \Phi(\mathbf{x},v,t)\end{psmallmatrix*}
		$
		has the same soln space as Newt's eqn.
		
		\textbf{Princ:} (Newt's 1$^{st}$ Law). A particle (or many-body system of particles) is freely-moving if $\Phi(\mathbf{x},\dot{\mathbf{x}},t)=0$. In this case, Newt's eqn iss $\ddot{\mathbf{x}}(t)=0$. If the initial state of such a freely moving system is $\mathbf{x}(0)=\mathbf{x}_0$, $\dot{\mathbf{x}}(0)=v_0$, then the corresponding Newt's eqn is solved by $\mathbf{x}(t)=\mathbf{x}_0+tv_0$. Thus the freely moving body moves on a straight worldline at constant vel..
		
		\textbf{Princ:} (Newt's 2$^{nd}$ Law). A force is a map $F:\mathbb{R}^3\to\mathbb{R}^3$ which affects the vel. of a particle. For a particle of (inertial) mass $m$, the Newt's eqn in a force field becomes
		$
		F(\mathbf{x})=m\ddot{\mathbf{x}},\text{or} \Phi(\mathbf{x},\dot{\mathbf{x}},t)=\frac{F(\mathbf{x})}{m}.
		$
		
		\textbf{Dfn:} A point $\mathbf{x}\in\mathbb{R}^3$ for which $F(\mathbf{x})=0$ is an \textbf{equil. point} of the force field $F$.
		
		\hrule
		\vspace*{-3pt}
		\subsubsection*{\tiny{2. Galilean Relativity}}
		\vspace*{-5pt}
		
		\textbf{Dfn:} A \textbf{Galilean transform} is an affine transform which preserves time intervals and the distances between simul. events.
		
		\textbf{Lem:} A Galil. transform is a linear transform of form
		$
		\begin{psmallmatrix*}1 & 0 & s\\ v & R & s\\ 0 & 0 & 1\end{psmallmatrix*}
		$
		with $s,v\in\mathbb{R}^3$, $R:\mathbb{R}^3\to\mathbb{R}^3$ an orthog transform.
		Consider three types: \\
		Type 1:
		$
		\begin{psmallmatrix*}1&0&s\\0&1&1\\0&0&1\end{psmallmatrix*}\begin{psmallmatrix*}t\\x\\1\end{psmallmatrix*}=\begin{psmallmatrix*}t+s\\x+s\\1\end{psmallmatrix*}
		$
		Type 2:
		$
		\begin{psmallmatrix*}1&0&0\\v&1&0\\0&0&1\end{psmallmatrix*}\begin{psmallmatrix*}t\\x\\1\end{psmallmatrix*}=\begin{psmallmatrix*}t\\x+vt\\1\end{psmallmatrix*}
		$
		Type 3:
		$
		\begin{psmallmatrix*}1&0&0\\0&R&0\\0&0&1\end{psmallmatrix*}\begin{psmallmatrix*}t\\x\\1\end{psmallmatrix*}=\begin{psmallmatrix*}t\\Rx\\1\end{psmallmatrix*}
		$
		Type 1 is translation/change of origin in $\mathbb{R}^4$. Type 2 is motion with fixed, uniform vel. $v$. Type 3 is re-orientation of spatial axes.
		
		\textbf{Princ:} (Galilean Relativity). The laws of motion (Newt's Eqn) are invariant under Galilean transforms - translation, fixed $\dot{x}$ motion and rotation of spatial axes.
		
		\hrule
		\vspace*{-3pt}
		\subsubsection*{\tiny{3. Conservation Laws}}
		\vspace*{-5pt}
		
		\textbf{Dfn:} A \textbf{conserved qtty} is a func of $(\mathbf{x},v)$ which remains constant in time - ie has time-deriv $0$ along solns of the eqns of motion.
		
		\textbf{Dfn:} $F:\mathbb{R}^3\to\mathbb{R}^3$ is \textbf{conservative} if it can be written as $F=-\nabla U$ for some func $U:\mathbb{R}^3\to\mathbb{R}$ where $\nabla$ is the gradient operator.
		
		\textbf{Dfn:} The \textbf{potential} of a conservative force field $F$ is $U(\mathbf{x})$. Potential is a func $U:\mathbb{R}^N\to\mathbb{R}$ s.t. $\ddot{\mathbf{x}}=-\nabla U$.
		
		\textbf{Dfn:} The \textbf{energy} of a particle of mass $m$ is
		$
		E(\mathbf{x},v)=\frac{m}{2}|v|^2+U(\mathbf{x})
		$
		where $T(v)=\frac{m}{2}|v|^2$ is kinetic energy and $U(\mathbf{x})$ is potential energy. $E=T+U$.
		
		\textbf{Lem:} In a conservative force field, the energy is a cnsrvd qtty.
		
		\textbf{Dfn:} Along some path $C\subset\mathbb{R}^3$ with endpoints $x_0$ and $x_1$, the \textbf{Work} $W$, done by a force field $F$ from $x_0$ to $x_1$ via the path $C$ is
		$
		W = \int_{x_0}^{x_1} F\cdot dx = \int_C F\cdot dx.
		$
		Paramise by $x(t)$ s.t. $x_0=x(t_0)$, $x_1=x(t_1)$ then
		$
		W = \int_{t_0}^{t_1} F\cdot vdt.
		$
		
		\textbf{Thm:} Let $F$ be a force field, then TFAE.
		\begin{enumerate}[leftmargin=*,topsep=0pt,itemsep=0pt]
			\item $F$ is conservative, $F=-\nabla U$ for some $U:\mathbb{R}^3\to\mathbb{R}$.
			\item The work of $F$ along any path $C$ connecting $x_0$ to $x_1$ depends only on $x_0,x_1$ and not on the shape of $C$, and if $C$ is a closed path then the work of $F$ along $C$ is $0$.
			\item The curl of $F$ vanishes, $\nabla\times F=0$.
		\end{enumerate}
		
		\textbf{Dfn:} Motion in 1D for a conservative force is governed by
		$
		m\ddot{x} = -\frac{dU}{dx}.
		$
		If $U(x)$ is linear then motion is the same as the example of Galil. gravity. When $U(x)$ is quadratic --- this is ``simple harmonic motion''.
		
		\textbf{Princ:} (Hooke's Law). The restoring force of a spring is well approxed by the potential
		$
		U(x)=\frac{k}{2}x^2
		$
		where $k>0$ is a constant characterising stiffness of the spring. The resulting force by Hooke's Law is $F=-kx$ and Newt's eqn becomes $m\ddot{x}=-kx$. The amplitude of the spring's oscillation is $A=\sqrt{2E/k}$. Physical trajectories in the state space $\mathbb{R}^2$.
		
		\textbf{Lem:} In any 1D conservative system,
		$
		E = \frac{m}{2}v^2+U(x),
		$
		is a constant of motion. As $\frac{m}{2}v^2\geq 0$, it follows that $E\geq U(x)$ with equality iff $v=0$. Configurations with $E<U(x)$ are not allowed.
		
		With the system at fixed energy $E$, there are two permitted regions of motion: for $x(t)\in[a,b]$ the motion will be oscillatory (in fact periodic); and for $x(t)\in[c,\infty)$ motion will tend towards $x\to\infty$. In the oscillatory case the period of motion is
		$
		T = \sqrt{2m}\int_a^b 1/\sqrt{E-U(x)} \text{ dx}.
		$
		This integral converges iff $a$ and $b$ are not CPs of the potential. If $b$ is such a CP the particle takes an `infinite' amount of time to reach $b$.
		
		Now suppose the energy is increased until it coincides with a local max of the potential as shown. If the particle starts from rest at $x(0)=a$, the previous formula suggests periodic motion between $a$ and $b$ but this can't happen now as $b$ is an equil.: if the particle reaches $b$ it must have $0$ vel. and $0$ force field hence the particle is fixed at $b$ forever. What really happens is the particle never reaches $b$. Consider the time it takes for the particle to get from $a$ to $b$
		$
		\Delta t(a\to b) = \sqrt{\frac{m}{2}}\lim_{\epsilon\to 0^+}\int_a^{b-\epsilon}\frac{dx}{\sqrt{E-U(x)}}.
		$
		
		Now Taylor expand around $x=b$:
		$
		U(x) = U(b)+U'(b)(x-b)+\tfrac{1}{2}U''(b)(x-b)^2+O((x-b)^3).
		$
		If $U'(b)=0$ and $E=U(b)$:
		$
		E-U(x) = -\tfrac{1}{2}U''(b)(b-x)^2+O((x-b)^3).
		$
		As $b$ is a local max of the potential, $U''(x)<0$ and so the integral is approx (ignoring the factor of $m/2$):
		$
		2|U''(b)|^{-1/2}\int\frac{dx}{b-x} \approx -2|U''(b)|^{-1/2}\log(b-x)
		$
		which is unbdd as $x\to b$. Even if $U''(b)=0$ a similar argument shows unbddness. Thus, if $b$ is a CP of the potential, the particle never reaches $b$.
		
		\hrule
		\vspace*{-3pt}
		\subsubsection*{\tiny{4. The Two-Body Problem}}
		\vspace*{-5pt}
		
		For $x_1,x_2\in\mathbb{R}^3$ the positions of the two particles, masses $m_1$ and $m_2$ respectively; let $U(x_1,x_2)=U(|x_1-x_2|)$ be the potential of their interaction.
		
		Newt's eqn for the two body system is:
		$
		m_1\ddot{x}_1=-\nabla_1 U, m_2\ddot{x}_2=-\nabla_2 U
		$
		where $\nabla_1$ stands for the gradient operator wrt $x_1$, etc. Since the potential only depends on the distance between the particles, $U=U(|x_1-x_2|)$ it follows that
		$
		\nabla_1 U = -\nabla_2 U, \text{hence} m_1\ddot{x}_1+m_2\ddot{x}_2=0.
		$
		This in turn implies that the qtty $p_c:=m_1\dot{x}_1+m_2\dot{x}_2$ is cnsrvd. This can be seen as the momentum of the center-of-mass (CoM). The centre-of-mass coordinate $x_c$ for the two-body system as:
		$
		x_c := \frac{m_1 x_1}{m_1+m_2}+\frac{m_2 x_2}{m_1+m_2}.
		$
		Then $p_c=(m_1+m_2)\dot{x}_c$, so the cnsrvd qtty $p_c$ is referred to as the centre-of-mass momentum of the two-body system.
		
		Rewrite the problem in terms of the relative coordinate $x:=x_1-x_2$ and $x_c$ s.t.:
		$
		\begin{psmallmatrix*}x\\x_c\end{psmallmatrix*}=\begin{psmallmatrix*}1 & -1\\ \frac{m_1}{m_1+m_2} & \frac{m_2}{m_1+m_2}\end{psmallmatrix*}\begin{psmallmatrix*}x_1\\x_2\end{psmallmatrix*}
		$
		$
		\begin{psmallmatrix*}x_1\\x_2\end{psmallmatrix*}=\begin{psmallmatrix*}\frac{m_2}{m_1+m_2} & 1\\ -\frac{m_1}{m_1+m_2} & 1\end{psmallmatrix*}\begin{psmallmatrix*}x\\x_c\end{psmallmatrix*}.
		$
		Newt's eqns in these new coordinates is $\ddot{x}_c=0$,
		$
		\mu\ddot{x}=-\nabla U(|x|),
		$
		where $\nabla$ denotes the gradient with respect to the relative coordinate $x$ and $\mu:=\frac{m_1 m_2}{m_1+m_2}$ is called the \textbf{reduced mass} of the two body system. These are decoupled, and thus the two body system is equivalent to a system of two non-interacting particles. A free particle located at $x_c$ of mass $m_1+m_2$, and a particle of mass $\mu$ located at $x$ and acting under the conservative force field $-\nabla U$. The 1$^{st}$ eqn of motion for the CoM can be solved as
		$
		x_c(t)=x_c(0)+\frac{p_ct}{m_1+m_2}.
		$
		The energy of the system then decomposes into the energies of each of the decoupled parts:
		$
		E=\frac{m_1}{2}|v_1|^2+\frac{m_2}{1}|v_2|^2+U(|x|)
		=\frac{m_1+m_2}{2}|\dot{x}_c|^2+\frac{\mu}{2}|\dot{x}|^2+U(|x|).
		$
		\vspace{1pt}
		\hrule
		\vspace*{-3pt}
		\subsubsection*{\tiny{5. Elastic Collisions}}
		\vspace*{-5pt}
		
		When $U(|x|)=0$ the two particles move freely, and $p_1+p_2=m_1v_1+m_2v_2$ is cnsrvd in addition to the energy. In this scenario an \textbf{elastic collision} is when the two free particles collide and emerge from this collision with potentially different momenta. Using Galil. transform (motion with a constant vel.) to set $v_2=0$ the conservation of momentum and energy then give the relations:
		$
		p_1 = p_1'+p_2',  \frac{m_1}{2}|v_1|^2=\frac{m_1}{2}|v_1'|^2+\frac{m_2}{2}|v_2'|^2.
		$
		Where $p_1',p_2',v_1'$ and $v_2'$ are the respective quantities after the collision. Thus the motion takes place in the plane spanned by $\{p_1',p_2'\}$. Breaking the conservation of momentum into components parallel and perpendicular to $p_1$ gives:
		$
		m_1|v_1| = m_1|v_1'|\cos\theta_1+m_1|v_2'|\cos\theta_2,
		$
		$
		0 = m_1|v_1'|\sin\theta_1+m_2|v_2'|\sin\theta_2,
		$
		where $\theta_1,\theta_2$ are the angles between $p_1',p_2'$ and $p_1$ respectively. (leaving 3 eqns to solve for 4 unknowns). s.t. can't model the collision, but can put constraints on the outgoing motion. This is done by finding constraints on $\theta_1$.
		
		This is easiest relative to the CoM frame, where the CoM vel. is
		$
		v_c = \frac{m_1}{m_1+m_2}v_1
		$
		is constant (as the particle motion is free). Then the relative velocities of each particle are $u_1:=v_1-v_c$ and $u_2:=-v_c$ before the collision then $u_1':=v_1'-v_c$ and $u_2':=v_2'-v_c$ after the collision. Working in the CoM frame motion must be co-linear (otherwise CoM vel. would be 0). The conservation of energy gives
		$
		\frac{m_1}{2}|u_1|^2+\frac{m_2}{2}|u_2|^2 = \frac{m_1}{2}|u_1'|^2+\frac{m_2}{2}|u_2'|^2.
		$
		Subtracting the CoM energy from both sides and using the conservation of momentum gives $u_2=-\frac{m_1}{m_2}u_1$ and $u_2'=-\frac{m_1}{m_2}u_1'$. Plugging this into the conservation of energy eqn gives $|u_1|=|u_1'|$ and so $|u_2|=|u_2'|$. So, relative to the CoM, the velocities are rotated by an angle, $\theta$, which is not constrained by the conservation of momentum or energy.
		
		Relate $\theta$ to $\theta_1$ by decomposing $v_1'=u_1'+v_c$ into components parallel and perpendicular to $v_c$:
		$
		|v_1'|\cos\theta_1 = |u_1'|\cos\theta+|v_c|,  |v_1'|\sin\theta_1=|u_1'|\sin\theta.
		$
		Which can be rearranged to give:
		$
		\tan\theta_1 = \frac{\sin\theta}{\cos\theta+(m_1/m_2)}.
		$
		
		\hrule
		\vspace*{-3pt}
		\subsubsection*{\tiny{5. Central Force Fields}}
		\vspace*{-4pt}
		
		\textbf{Dfn:} A \textbf{central force} is one that points in the direction of the line connecting the particle on which it acts with the origin for all configs.
		
		\textbf{Dfn:} (Central Force Field). A force field $F:\mathbb{R}^3\to\mathbb{R}^3$ is called \textbf{central} if it takes the form $F(\mathbf{x})=f(\mathbf{x})\mathbf{x}$, for some func $f:\mathbb{R}^3\to\mathbb{R}$.
		
		\textbf{Thm:} (5.1). Let $F:\mathbb{R}^3\to\mathbb{R}$ be a conservative force field. Then\\
			$F$ is central $\iff$ $F(\mathbf{x})=f(|F|)\mathbf{x}$ $\iff$ $F(\mathbf{x})=-\nabla U(|\mathbf{x}|)$.
		
		\textbf{Dfn:} For a particle of mass $m$ the \textbf{ang mom} is
		$
		L=\begin{psmallmatrix*}L_1\\L_2\\L_3\end{psmallmatrix*}=\begin{psmallmatrix*}x_2p_3-x_3p_2\\ x_3p_1-x_1p_3\\ x_1p_2-x_2p_1\end{psmallmatrix*}.
		$
		aka $L:=\mathbf{x}\times p$.
		The components of the ang mom have dim $[L]=ML^2T^{-1}$.
		
		\textbf{Thm:} (5.2). In any central force field, ang mom is a cnsrvd qtty.
		
		\textbf{Cor:} (5.1). Motion in any central force field is restricted to a plane in $\mathbb{R}^3$.
		
		 Some Conditions for only bound motion: $\lim_{r\to\infty}U_{\mathrm{eff}}(r)=\infty$. Furthermore, by showing $\lim_{r\to 0^+}U_{\mathrm{eff}}=\infty$ and $U_{\mathrm{eff}}$ is finite for all $r\in(0,\infty)$ gives us that the motion can stay $r\in(0,\infty)$, rather than getting stuck at $r=0$. Conditions for only scattering motion: $\lim_{t\to\infty}U_{\mathrm{eff}}$ is finite, no minima at finite $r$, something like a wall at $r=0$, e.g.\ $\lim_{r\to 0^+}U_{\mathrm{eff}}=\infty$.
		
		\textbf{1D Effective Description}
		
		Use a Galil. rotation to orient our axes so that motion in a central field takes place in the $(x_1,x_2)$-plane:
		$L=(0, 0, m\ell)^T$
		Working in polar coords $x_1=r\cos\theta$, $x_2=r\sin\theta$ for the plane of motion. It follows momentum of a body subject to a central force is:
		$
		p = m\dot{\mathbf{x}} = m\begin{psmallmatrix*}\dot{r}\cos\theta - r\dot\theta\sin\theta\\ \dot{r}\sin\theta+r\dot\theta\cos\theta\\ 0\end{psmallmatrix*}.
		$
		Conservation of ang mom then implies:
		$
		L_3 = x_1p_2-x_2p_1 = mr^2\dot\theta\text{ is a constant and in particular}
		$
		$
		\ell = r^2\dot\theta\text{ is a cnsrvd qtty.}
		$
		
		\textbf{Princ:} (Kepler's Area Law). The areal vel. of a body moving in a central force field is constant.
		
		Adapting motion in a central field in terms of polar coords gives:
		$
		\frac{dp}{dt}=m\begin{psmallmatrix*}(\ddot{r}-\dot\theta^2 r)\cos\theta - (\ddot\theta r+2\dot\theta\dot{r})\sin\theta\\
			(\ddot{r}-\dot\theta^2 r)\sin\theta + (\ddot\theta r+2\dot\theta\dot{r})\cos\theta\\ 0\end{psmallmatrix*}.
		$
		Using conservation of ang mom, taking a time deriv of $\ell=r^2\dot\theta$ gives the identity $2\dot{r}\dot\theta+r\ddot\theta=0$ which simplifies the motion to
		$
		\frac{dp}{dt}=m\!\left(\ddot{r}-\frac{\ell^2}{r^3}\right)(\cos \theta, \sin \theta, 0)^T
		$
		By Newt's 2$^{nd}$ law $\dot{p}=F$ while Thm 5.1 implies
		$
		F(\mathbf{x})=-\frac{U'(|\mathbf{x}|)}{|\mathbf{x}|}\mathbf{x}=-U'(r)(\cos\theta, \sin\theta, 0)^T.
		$
		So Newt's eqn for motion in a conservative central force field reduces to a 1D eqn:
		$
		m\ddot{r}=-U'(r)+\frac{m\ell^2}{r^3}\text{ or } m\ddot{r}=-U'_{\mathrm{eff}}(r)
		$
		where
		$
		U_{\mathrm{eff}}:=U(r)+\frac{m\ell^2}{2r^2}, U_{\mathrm{eff}}(|\mathbf{x}|)=U(|\mathbf{x}|)+\frac{|L|^2}{2m|\mathbf{x}|^2}.
		$
		Conservation of energy for the 1D effective description of motion states $E=m\dot{r}^2/2+U(r)+\frac{m\ell^2}{2r^2}$ is cnsrvd. Note that $\frac{m\ell^2}{2r^2}=mr^2\dot\theta^2/2$ for centrifugal energy. NOTE that $E=U_{\mathrm{eff}}$ are NOT CPs with 0 vel.; just local mins/max of $r(t)$, there can still be movement in the angular direction (ie possible for $\dot\theta \ne 0$). The energy eqn above is equal to the total energy of the original 2D eqn as $m|\dot{\mathbf{x}}|^2=m\dot{r}^2/2+mr^2\dot\theta^2$.
		
		In the 1D effective description, there are two classes of motion: $r\in[r_{\min},\infty)$ (scattering), $r\in[r_{\min},r_{\max}]$ (bound). In scattering, the body approaches the origin from some initial (possibly infinite) radius, achieves a closest approach of $r_{\min}$ ($E=U_{\mathrm{eff}}$) and then flies away to $\infty$. In bound, the body orbits the origin between the closest and furthest radii i.e. bound orbits lie in the annulus bdd by $[r_{\min},r_{\max}]$ in the plane. By conservation of energy and ang mom
		$
		\frac{d\theta}{dr}=\frac{\dot\theta}{\dot{r}}=\pm \ell \left(r^2\sqrt{(2/m)(E-U(r))-\ell^2/(r^2)}\right)^{-1}.
		$
		As here, Newt's eqs are invariant under time inversion, the angle through which $\mathbf{x}(t)$ passes as it goes from $r_{\min}$ to $r_{\max}$ is half of $\Delta\theta$ so
		$
		\Delta\theta = 2\ell\int_{r_{\min}}^{r_{\max}}\frac{dr}{r^2\sqrt{(2/m)(E-U(r))-\frac{\ell^2}{r^2}}}
		=2\ell\int_{r_{\min}}^{r_{\max}}\frac{dr}{r^2\sqrt{(2/m)(E-U_{\mathrm{eff}}(r))}}.
		$
		
		\textbf{Lem:} (5.1). The orbits of bound motion in a conservative central force is closed iff $\Delta\theta=2\pi n/m$ where $n,m\in\mathbb{Z}$ are coprime.
		
		\hrule
		\vspace{-3pt}
		\subsubsection*{\tiny{6. The Kepler Problem}}
		\vspace*{-4pt}
		
		\textbf{Dfn:} The \textbf{Kepler problem} is the general problem of motion in conservative central forces with potentials of form
		$
		U(r) = -k/r,
		$
		where $k$ is a physical constant, dim = $[k]=ML^3T^{-2}$. $k>0 \Rightarrow$attractive potential, $k<0 \Rightarrow$repulsive.

		The 1D effective potential for the Kepler problem is
		$
		U_{\mathrm{eff}}(r) = -\frac{k}{r}+\frac{m\ell^2}{2r^2},
		$
		where $\ell=r^2\dot\theta$ is a constant by the conservation of ang mom.
		
		This maps onto the 2-body problem by decoupling the (free) CoM motion and taking $r=|x_1-x_2|$, $m=\mu$ (the reduced mass of the two-body problem).
		
		When $k<0$ and the potential is repulsive, $U_{\mathrm{eff}}(r)>0$, $U'_{\mathrm{eff}}(r)<0$ for all $r$, so only scattering motion is possible, with the two bodies reaching a minimal separation of $r_{\min}$ determined by $E=U_{\mathrm{eff}}(r_{\min})$.
		
		When $k>0$ the potential is attractive and bound motion is possible for configurations where $E<0$. In this case the radial motion is bdd by the roots of $E=U_{\mathrm{eff}}(r)$; this is a quadratic eqn in $r^{-1}$, giving
		$
		\frac{1}{r_{\min}}=\frac{k}{m\ell^2}\left(1+\sqrt{1-\frac{2|E|m\ell^2}{k^2}}\right)
		$
		$
		\frac{1}{r_{\max}}=\frac{k}{m\ell^2}\left(1-\sqrt{1-\frac{2|E|m\ell^2}{k^2}}\right).
		$
		
		\textbf{Lem:} For $k>0$, $E<0$ and $\ell\neq 0$, the turning angle for bound orbits of the Kepler problem is $\Delta\theta=2\pi$.
		
		The min of the effective potential with $k>0$ corresponds to the single (finite) root of
		$
		U'_{\mathrm{eff}}(r)=k/(r^2)-m\ell^2/(r^3)=0\implies r=m\ell^2/k.
		$
		The value of the effective potential at this radius is
		$
		U^{\min}_{\mathrm{eff}}=-\frac{k^2}{2m\ell^2}.
		$
		When $E=U^{\min}_{\mathrm{eff}}$, $r_{\min}=r_{\max}=m\ell^2/k$ so the bound orbit is circular. Then by conservation of ang mom,
		$
		\dot\theta = \frac{k^2}{m^2\ell^3},
		$
		and the period of the orbit is
		$
		T=\frac{2\pi m^2\ell^3}{k^2}.
		$
		
		\textbf{Dfn:} (Laplace-Runge-Lenz vector).
		$
		A := \dot{\mathbf{x}}\times L + U\mathbf{x}.
		$
		\hrule
		\vspace*{-3pt}
		\subsubsection*{\tiny{7. Kepler's Laws of Planetary Motion}}
		\vspace*{-5pt}
		
		The Laplace-Runge-Lenz (L-R-L) vector is orthogonal to $L$ ($A\cdot L=0$), thus $A$ always lies in the plane of motion for the Kepler problem.
		
		\textbf{Dfn:} Define $\varphi$ as the angle between $\mathbf{x}$ and $A$ in the plane of motion, then
		$
		\mathbf{x}\cdot A = |\mathbf{x}||A|\cos\varphi.
		$
		By the dfn of the Laplace-Runge-Lenz vec:
		$
		\mathbf{x}\cdot A = \mathbf{x}\cdot(\dot{\mathbf{x}}\times L)+U|\mathbf{x}|^2,
		$
		and taking $U=-k/|\mathbf{x}|$ and using the identity $a\cdot(b\times c)=c\cdot(a\times b)$ gives
		$
		\mathbf{x}\cdot A = m\ell^2-k|\mathbf{x}|.
		$
		Equating with the earlier expression gives
		$
		|\mathbf{x}|\left(1+\frac{|A|}{k}\cos\varphi\right)=\frac{m\ell^2}{k},
		$
		paramizing the radial motion $r=|\mathbf{x}|$ in terms of the angle $\varphi$ between the position and L-R-L vectors.
		
		\textbf{Dfn:} A \textbf{conic section} in $\mathbb{R}^3$ is a planar curve obtained by intersecting an affine plane with a right circular double cone.
		
		If $(r,\varphi)$ are polar coordinates on the plane, then the curve of a conic section is given by
		$
		r(1+\varepsilon\cos\varphi)=\lambda,
		$
		where the focus of the conic section is at $r=0$, the param $\varepsilon\geq 0$ is called the \textbf{eccentricity} and $\lambda>0$ is the \textbf{semi-latus rectum}. The latter is the distance from the focus at $r=0$ to the conic section in the direction of the perpendicular to the major axis.
		
		\textbf{Dfn:} The named cases of conic sections are given by the eccentricity:\\
		Circle when $\varepsilon=0$, Ellipse when $\varepsilon<1$
		(closed curves); 
		Parabola when $\varepsilon=1$, Hyperbola when $\varepsilon>1$
		(open curves).
		
		\textbf{Lem:} Motion in the Kepler problem is given by conic sections in the plane perpendicular to $L$ with eccentricity $\varepsilon=\frac{|A|}{k}$ and semi-latus rectum
		$
		\lambda=\frac{m\ell^2}{k},
		$
		
		\textbf{Lem:} (7.2). For the attractive Kepler problem ($k>0$), $E<0$ iff the eccentricity of the resulting conic section is $\varepsilon<1$.
		
		\textbf{Princ:} Kepler's $1^{st}$ Law: Bounded orbits in the Kepler problem are ellipses with the origin as one of the foci. (The orbit of a planet is an ellipse with the Sun at one of the two foci.)
		
		\textbf{Princ:} Kepler's $2^{nd}$ Law: The areal vel. of motion (change in area w.r.t. time) in the Kepler problem is constant.
		
		\textbf{Princ:} Kepler's $3^{rd}$ Law: The period of a bdd orbit in the Kepler problem is proportional to the $\frac{3}{2}$ power of the length of the semi-major axis of the corresponding ellipse.
		
		The min and max distances between the two bodies can be characterised by the quantities $a$ and $b$; the lengths of the semi-major ($a$) and semi-minor ($b$) axes of the ellipse.
		$
		\mathrm{Area(ellipse)}=\int_0^T\frac{dA}{dt}(t)dt,
		$
		where $A(t)$ is the area swept out by the trajectory $(r(t),\theta(t))$ as a func of time. By Kepler's 2$^{nd}$ law
		$
		\frac{dA}{dt}(t) = \frac{r^2(t)}{2}\dot\theta(t) = \ell
		$
		is constant so
		$
		\mathrm{Area(ellipse)}=\int_0^T\ell/2dt=\ell T/2,
		$
		or equivalently $T=2\mathrm{Area(ellipse)}/\ell$.
		
		Note $2a=r_{\min}+r_{\max}$. In the polar description of an ellipse $(\varepsilon<1)$, $r_{\min}$ occurs at $\varphi=0$ and $r_{\max}$ occurs at $\varphi=\pi$, so 
		$
		r_{\min}=\frac{\lambda}{1+\varepsilon}\text{and} r_{\max}=\frac{\lambda}{1-\varepsilon}.
		$
		Thus the length of the semi-major axis relates to the eccentricity and semi-latus rectum by
		$
		a=\frac{\lambda}{1-\varepsilon^2},\Leftrightarrow r(1+\varepsilon\cos\varphi)=a(1-\varepsilon^2).
		$
		Likewise the length of the semi-minor axis $b$ corresponds to the max value taken by $r\sin\varphi = \frac{a(1-\varepsilon^2)\sin\varphi}{1+\varepsilon\cos\varphi}.
		$
		Thus the length of the semi-minor axis is
		$
		b=\left.\frac{a(1-\varepsilon^2)\sin\varphi}{1+\varepsilon\cos\varphi}\right|_{\max}=a\sqrt{1-\varepsilon^2}.
		$
		Now express $a$ and $b$ in terms of the physical params of the bound orbit for
		$
		a=k/(2\left|E\right|), b=\sqrt{(m\ell^2)/(2\left|E\right|)}.
		$
		
		
		\textbf{Lem:} (7.3). The area of an ellipse with semi-major and semi-minor axes of lengths $a$ and $b$, respectively, is $\mathrm{Area(ellipse)}=\pi ab$.
		
		Substituting into $T$ gives $T=\frac{2\pi ab}{\ell}=\frac{2\pi k\sqrt{m}}{(2|E|)^{3/2}}$, and our identity for $b$ gives
		$
		T=2\pi a^{3/2}\sqrt{m/k},
		$
		in accordance with Kepler's 3$^{rd}$ Law.
		
		\hrule
		\vspace*{-3pt}
		\subsubsection*{\tiny{8. Scattering in Central Force Fields}}
		\vspace*{-5pt}
		
		Parabolic orbits ($\varepsilon=0$) are only possible in the attractive ($k>0$) Kepler problem, for configurations with $E=0$. The min radius $r_{\min}$ relative to the focus at the origin corresponds to $U_{\mathrm{eff}}(r_{\min})=0$; so
		$
		r_{\min}^{\mathrm{para}}=\frac{m\ell^2}{2k}.
		$
		WLOG paramise time using a Galil. transform s.t. $r(0)=r_{\min}$, which means that $\varphi(0)=0$. Then
		$
		t=\frac{m^2\ell^3}{4k^2}\left[\tan\frac{\varphi}{2}+\frac{1}{3}\tan^3\!\left(\frac{\varphi}{2}\right)\right].
		$
		
		\textbf{Dfn:} For the cnsrvd quantities $(E,\ell)$ and geometric params $(\varepsilon,\lambda)$ of a conic section:
		$
		\varepsilon=\sqrt{1+\frac{2m\ell^2}{k^2}E},\lambda=\frac{m\ell^2}{k},
		$
		$
		E=\frac{k(\varepsilon^2-1)}{2\lambda},\ell=\sqrt{\frac{k\lambda}{m}}.
		$
		
		Hyperbolic orbits, for both attractive and repulsive potentials of the Kepler problem, correspond to $E>0$ configurations. As $\lim_{r\to\infty}U_{\mathrm{eff}}(r)=0$;
		
		\textbf{Dfn:} The \textbf{asymptotic speed}, $v_\infty$, is the positive real number s.t. $E=\frac{m}{2}v_\infty^2$ (since all energy at $r\to\infty$ has become kinetic).
		
		\textbf{Dfn:} The \textbf{impact param} $b$ associated with a given hyperbola is the min distance between either of its linear asymptotes and the origin.
		
		Now add two more dictionaries. Firstly, between $(E,\ell)$ and $(v_\infty,b)$:
		$
		v_\infty=\sqrt{\frac{2E}{M}}, b=\sqrt{\frac{m\ell^2}{2E}},
		$
		$
		E=\frac{M}{2}v_\infty^2,\ell=bv_\infty,
		$
		.Secondly, between $(\varepsilon,\ell)$ and $(v_\infty,b)$:
		$
		\varepsilon=\sqrt{1+\frac{m^2b^2v_\infty^4}{k^2}},\lambda=\frac{mb^2v_\infty^2}{k},
		$
		$
		v_\infty=\sqrt{\frac{k}{\lambda m}(\varepsilon^2-1)}, b=\frac{\lambda}{\sqrt{\varepsilon^2-1}}.
		$
		
		\textbf{Dfn:} In a hyperbol. orbit,  the closest $r_{\min}$ to the origin is
		$
		r_{\min}=-\frac{k}{mv_\infty^2}+\sqrt{b^2+\frac{k^2}{m^2v_\infty^4}}
		$
		
		As the body moves from $\infty$ to $r_{\min}$, the orbit passes through a turning angle $\chi$ before going back out to $\infty$. The \textbf{scattering angle} $\psi$ is the angle of deflection the orbit experiences between its initial and final asymptotic trajectories; it is related to the turning angle by
		$
		\psi=|\pi-2\chi|.
		$
		
		\textbf{Dfn:} $\chi$ can be written
		$
		\chi=\ell\int_{r_{\min}}^{\infty}\frac{dr}{r^2\sqrt{\frac{2}{m}(E-U(r))-\frac{\ell^2}{r^2}}}.
		$
		This can be substituted into $\psi=|\pi-2\chi|$ for the scattering angle.
		
		\textbf{Lem:} For hyperbol. scattering orbits in the Kepler problem, the scattering angle is
		$
		\psi = 2\arcsin\!\left(1/\left(\sqrt{1+\frac{m^2b^2v_\infty^4}{k^2}}\right)\right).
		$
		
		\hrule
		\vspace*{-3pt}
		\subsubsection*{\tiny{9. Small Oscillations and Stable Equilibria}}
		\vspace*{-5pt}
		
		Consider in 1D where Newt's eqn is $m\ddot{x}=-U'(x)$. If $U$ has a CP use a Galil. translation to set the CP to be at $x=0$, so that $U'(0)=0$. Along with the assumption that $U(0)=0$ (since only defined to a constant) means the Taylor expansion around the CP is
		$
		U(x) = \frac{U''(0)}{2}x^2+O(x^3).
		$
		Let $k:=U''(0)$, then if $k=0$ the cp is degen., stability is indeterminate. If $k>0$ then the cp is a stable (local min) and if $k<0$ it is unstable (local max).
		
		\textbf{Princ:} (9.1). Simple harmonic motion, or Hooke's law, provides a universal 1$^{st}$-approx to motion near any stable equil. in a 1D mechanical system.
		
		\textbf{Dfn:} (Radial Evolution Eq). For small perturbs for circular orbit, $r=r_0+\rho$ for $\rho\ll 1$ --- the radial evolution eqn is
		$
		m\ddot{\rho}=-U''_{\mathrm{eff}}(r_0)\rho.
		$
		This corresponds to SHM with freq $\omega=\sqrt{U''_{\mathrm{eff}}(r_0)/m}$.
		
		\textbf{Dfn:} The \textbf{Hessian matrix} $H$ is defined entry-wise by
		$
		H_{ij}(\mathbf{x})=\frac{\partial^2 U}{\partial q_i\partial q_j}(\mathbf{x}).
		$
		
		\textbf{Thm:} (Spectral thm for symmetric matrices) Let $H$ be a real-valued $3\times 3$ symmetric matrix with eigvals $k_i\in\mathbb{R}$ (for $i=1,2,3$) --- not necessarily distinct --- and corresponding non-zero, lin. ind. eigvecs $v_i\in\mathbb{R}^3$, obeying $Hv_i=k_iv_i$. After normalising the eigvecs so that $|v_i|=1\forall i$ and $v_i\cdot v_j=0$ for $i\neq j$, the matrix $S=[v_1v_2v_3]$ is orthogonal ($S^TS=I$) and
		$
		HS=SD, D=\mathrm{diag}(k_1,k_2,k_3).
		$
		
		\textbf{Thm:} (9.1). Motion near any stable equil. point is equivalent to three decoupled harmonic oscillators with eqns
		$
		m\ddot{y}=-Dy\Leftrightarrow\ddot{y}_i=-\frac{k_i}{m}y_i, i=1,2,3.
		$
		Specifically, the resulting solns are
		$
		y_i(t)=A_i\sin(\omega_i t+\varphi_i),\omega_i^2:=\frac{k_i}{m}
		$
		with $A_i$ and $\omega_i$ fixed as constants. The physical parametrisation of the mechanical system is obtained by $\mathbf{x}(t)=Sy(t)$.
		
		\textbf{Dfn:} The variables $y_i(t)$ describing oscillatory motion around a stable equil. are called the \textbf{normal modes} of the mechanical system, and the frequencies $\omega_i$ are called the \textbf{characteristic frequencies}.
		
		\hrule
		\vspace*{-3pt}
		\subsubsection*{\tiny{10. Near-Equilibrium Dynamics}}
		\vspace*{-5pt}
		
		Suppose $q_0\in\mathbb{R}^N$ is an equil. point of a general mechanical system whose configuration space is $N$-dim.al and is conservative (with a well-defined potential $U:\mathbb{R}^N\to\mathbb{R}$). Let $\{e_i\}_{i=1,\ldots,N}$ be the usual orthonormal basis of $\mathbb{R}^N$, so that a generic vector $q\in\mathbb{R}^N$ can be expanded as $q=\sum_{i=1}^n q_i e_i$, where $(q_1,\ldots,q_n)$ are the corresponding coordinates on the configuration space $\mathbb{R}^N$.
		
		Thus $\nabla U(q_0)=0$, and since potential is only defined up to constants, can set $U(q_0)=0$, in which case Taylor expanding the potential around this equil. gives
		$
		U(q)=\tfrac{1}{2}(q-q_0)\cdot H(q_0)(q-q_0)+O(|q-q_0|^3).
		$
		
		\textbf{Dfn:} Assume the kinetic energy of the system is quadratic in the velocities, and can be written in terms of a $N\times N$ symmetric, positive-definite mass matrix $M$, as:
		$
		T=\frac{1}{2}\dot{x}\cdot M\dot{x}=\frac{1}{2}\sum_{i,j=1}^N M_{ij}\dot{x}_i\dot{x}_j.
		$
		
		\textbf{Princ:} (10.1, Near-Equilibrium Normal Form). The dynamics of small displacements from equil. can be studied by:
		\begin{enumerate}[leftmargin=*,topsep=0pt,itemsep=0pt]
			\item diagonalizing and normalizing the mass matrix, then
			\item diagonalizing the potential and solving for the normal modes and characteristic frequencies of the system.
		\end{enumerate}
		
		\textbf{General Form of Solutions.} Let $M$ be the mass matrix. Then by spectral thm there exists an orthonormal matrix $S_1$ (made out of the normalised eigenvectors of $M$) s.t. $M'=S_1^T MS_1$, where $M'$ is diagonal with positive entries ($M$ is positive definite). Let $D_1$ be s.t.\ $M'=D_1^2$, then $M=S_1D_1^2S_1^T=(S_1D_1)(S_1D_1)^T$.
		
		Now, define $z:=(S_1D_1)^Tx=D_1S_1^Tx$, which is invertible as $x=S_1D_1^{-1}z$. In terms of $z$ the kinetic energy is
		$
		T=\frac{1}{2}\sum_{i=1}^N\dot{z}_i^2=\frac{1}{2}|\dot{z}|^2.
		$
		Thus normalising the near equil. dynamics. In these new variables the potential energy is given by
		$
		U=\frac{1}{2}z\cdot Kz, K=D_1^{-1}S_1^T HS_1D_1^{-1},
		$
		to leading order around the point of equil.. $K$ is symmetric since $H$ and $D_1$ are symmetric and $S_1$ is orthogonal. Thus by the spectral thm for symmetric matrices, there must exist another orthogonal matrix $S_2$ s.t. $D:=S_2^T KS_2$ is diagonal. Define $y:=S_2^Tz$, under which the form of kinetic energy is preserved
		$
		T=\frac{1}{2}|S_2\dot{y}|^2=\frac{1}{2}|\dot{y}|^2,
		$
		since $S_2$ is orthogonal, while the potential energy becomes
		$
		U=\frac{1}{2}y\cdot Dy.
		$
		In particular, the resulting linear eqns of motion of the system are $\ddot{y}=-Dy$, which are $N$-decoupled linear eqns of motion as $D$ is diagonal. Let $D=\mathrm{diag}(\lambda_1,\ldots,\lambda_N)$, then the decoupled eqns of motion are
		$
		\ddot{y}_i=-\lambda_i y_i, i=1,\ldots,N.
		$
		Each one falls into three cases:
		\begin{enumerate}[leftmargin=*,topsep=0pt,itemsep=0pt]
			\item $\lambda_i>0$: oscillatory solns with characteristic frequency $\omega_i=\sqrt{\lambda_i}$ and normal mode $y_i(t)=A_i\sin(\omega_i t+\varphi_i)$.
			\item $\lambda_i=0$: linear, or zero-mode, solns with characteristic frequency $\omega_i=0$ and normal mode $y_i(t)=a_i+b_it$.
			\item $\lambda_i<0$: exponential solns with imaginary characteristic frequency and normal mode $y_i(t)=A_i\exp\!\left(\sqrt{|\lambda_i|}t\right)+B_i\exp\!\left(-\sqrt{|\lambda_i|}t\right)$.
		\end{enumerate}
		
		\textbf{Dfn:} (10.1, Equilibrium stability). A point of equil. is:
		\begin{itemize}[leftmargin=*,topsep=0pt,itemsep=0pt]
			\item \textbf{Stable} if $\lambda_i>0$ for all $i=1,\ldots,N$,
			\item \textbf{Semi-stable} if $\lambda_i\geq 0$ for all $i=1,\ldots,N$ and
			\item \textbf{Unstable} if there is at least one $i\in\{1,\ldots,N\}$ for which $\lambda_i<0$.
		\end{itemize}
		
		\hrule
		\vspace*{-3pt}
		\subsubsection*{\tiny{11. Variational Calculus}}
		\vspace*{-5pt}
		
		\textbf{Dfn:} (Functional). A functional is a map from a set of funcs to $\mathbb{R}$.
		
		In particular, consider $k$-times cont.ly diffable funcs mapping from some closed domain $\Omega\subseteq\mathbb{R}^m$ to $\mathbb{R}^n$; this set of funcs is $C^k(\Omega,\mathbb{R}^n)$. Consider functionals $S:C^k(\Omega,\mathbb{R}^n)\to\mathbb{R}$ of form
		$
		S[f]=\int_\Omega L(f(x),\nabla f(x),\ldots,f^{(k)}(x);x)d\mu_\Omega,
		$
		where $f\in C^k(\Omega,\mathbb{R}^n)$, $L:\mathbb{R}^{m+n\frac{m^{k+1}-1}{m-1}}\to\mathbb{R}$ is a smooth, real valued func of its arguments and $d\mu_\Omega$ is an integral measure on $\Omega$.
		
		\textbf{Dfn:} (Trajectories) Let $x:[0,1]\to\mathbb{R}^n$ be a trajectory with fixed boundary conditions $x(0)=P$, $x(1)=Q$ for 2 points $P,Q\in\mathbb{R}^n$. $C_{P,Q}$ is the space of all such trajectories.
		
		\textbf{Dfn:} (endpoint fixed variations). Let $x(t),y(t)$ be any two paths in $C_{P,Q}$, and define $\varepsilon(t):=x(t)-y(t)$. As such, $\varepsilon:[0,1]\to\mathbb{R}^2$, $\varepsilon(0)=\varepsilon(1)=0$.
		
		Such a $\varepsilon(t)$ is called an \textbf{endpoint fixed variation}: it represents a shift in the initial path $x(t)$ which preserves the endpoints.
		
		A trajectory $x$ is an extremum of a functional $S[x]$ iff
		$
		\left.\frac{d}{ds}S[x+s\varepsilon]\right|_{s=0}=0,
		$
		for all endpoint fixed variations $\varepsilon:[0,1]\to\mathbb{R}^n$ with $\varepsilon(0)=\varepsilon(1)=0$.
		
		\textbf{Thm:} (Fundamental lemma of variational calculus). Let $f:[0,1]\to\mathbb{R}^n$ be a cont. func with
		$
		\int_0^1 f(t)\cdot h(t)=0,
		$
		for all $h\in C^\infty([0,1],\mathbb{R}^n)$ with $h(0)=h(1)=0$. Then $f(t)\equiv 0$.
		
		\hrule
		\vspace*{-3pt}
		\subsubsection*{\tiny{12. Euler-Lagrange Eqns}}
		\vspace*{-5pt}
		
		\textbf{Dfn:} A \textbf{Lagrangian} is a smooth func $L:\mathbb{R}^{2n+1}\to\mathbb{R}$, sending $(x,\dot{x},t)\mapsto L(x,\dot{x};t)$. The \textbf{action} associated to $L$ is the functional
		$
		S[x]=\int_0^1 L(x,\dot{x};t)
		$
		which is a map $S:C_{P,Q}\to\mathbb{R}$.
		
		Consider the Lagrangian given by $L=|\dot{x}(t)|$. The action associated to this Lagrangian is the arclength of the path $x$ from $P$ to $Q$.
		
		\textbf{Thm:} (12.1). A curve $x:[0,1]\to\mathbb{R}^n$ is an extremum of the action functional iff the E-L eqns
		$
		\frac{\partial L}{\partial x}=\frac{d}{dt}\frac{\partial L}{\partial\dot{x}}\Leftrightarrow\frac{\partial L}{\partial x_i}=\frac{d}{dt}\frac{\partial L}{\partial\dot{x}_i}.
		$
		Useful from Proof of Thm 12.1:
		$
		\frac{d}{ds}S[x+s\varepsilon]\big|_{s=0}=\int_0^1\frac{d}{ds}L(x,\dot{x};t)\big|_{s=0}dt
		=\int_0^1\!\!\left(\frac{\partial L}{\partial x}\cdot\varepsilon+\frac{\partial L}{\partial\dot{x}}\cdot\dot\varepsilon\right)dt.
		$
		
		\textbf{Lem:} (12.1). If $\frac{\partial L}{\partial x_i}=0$ for some $i\in 1,\ldots,N$, then the generalised momentum component $p_i$ is a cnsrvd qtty along curves obeying the E-L eqns.
		
		\textbf{Prop:} (12.1). If the Lagrangian does not explicitly depend on time, $\frac{\partial L}{\partial t}=0$, then the energy is a cnsrvd qtty along curves obeying the E-L eqns.
		
		\textbf{Princ:} (Hypersurfaces). Generalise the framework of extremising actions with fixed B.C.s, and let $x(0),x(1)$ be constrained to lie on hypersurfaces in $\mathbb{R}^n$, rather than points. Define the hypersurfaces $H_0,H_1\subset\mathbb{R}^n$ by
		$
		H_0=\{x\in\mathbb{R}^n|g_0(x)=0\}, H_1=\{x\in\mathbb{R}^n|g_1(x)=0\}
		$
		where $g_0,g_1:\mathbb{R}^n\to\mathbb{R}$ are smooth funcs. s.t. extremise the action subject only to $x(0)\in H_0$ and $x(1)\in H_1$.
		
		\textbf{Thm:} Let $x:[0,1]\to\mathbb{R}^n$ be subject to $x(0)\in H_0$ and $x(1)\in H_1$ for hypersurfaces $H_0,H_1\subset\mathbb{R}^n$. Then $x$ is an extremum of the action iff the E-L eqns are satisfied and the generalised momentum $p(t)$ obeys: $p(0)$ is normal to $H_0$ at $x(0)$ and $p(1)$ is normal to $H_1$ at $x(1)$.
		
		\textbf{Ex:} Using the arclength functional find the shortest path from $(2,1)$ ending on the hyperbola $xy=1$. SOL: By earlier examples on the arclength functional, the soln must be a straight line. Furthermore by the theory of general variations these extrema must be sols to the EL eqns and normal to the boundary hyper-surface at the point of intersection. So seek a line of form $x(t)=(1-t)\binom{2}{1}+t\binom{x_1}{y_1}$, which is normal to the parabola at $(x_1,y_1)$. Here $g(x,y)=xy-1$, thus $\nabla g(x_1,y_1)=\binom{y_1}{x_1}$. Also, $\dot{x}=\binom{x_1-2}{y_1-1}$. Thus the soln being normal to the hyperbola is equivalent to the determinant of the matrix of $\nabla g(x_1,y_1)$ and $\dot{x}$ being equal to zero. Solve the resulting eqn using the relations in $g$ to give the solns for $x_1$ and $y_1$.
		
		\hrule
		\vspace*{-3pt}
		\subsubsection*{\tiny{13. The 2nd Variation}}
		\vspace*{-5pt}
		
		\textbf{Dfn:} (Symmetric bi-linear forms and quadratic forms). A symmetric bi-linear form on $\mathbb{R}^n$ is a map $B:\mathbb{R}^n\times\mathbb{R}^n\to\mathbb{R}$ which is linear in each of its arguments and symmetric. In other words $\forall x,y\in\mathbb{R}^n$ and $\forall a,b\in\mathbb{R}$
		$
		B(ax,by)=abB(x,y), B(x,y)=B(y,x).
		$
		The bi-linear property means that such a $B$ can be represented as an $n\times n$ matrix, while the symm property means this matrix is symmetric.
		
		\textbf{Dfn:} (Positive/negative-definite). A symmetric bi-linear form is positive(negative)-definite if it obeys $B(x,x)>0$ ($B(x,x)<0$) for all $x\in\mathbb{R}^n\setminus\{0\}$.
		
		\textbf{Lem:} (13.1). A symmetric bi-linear form is positive(negative)-definite iff all of its eigenvalues are positive (negative).
		
		\textbf{Dfn:} For any symmetric bi-linear form $B$ there is an associated quadratic form $Q:\mathbb{R}^n\to\mathbb{R}$, $Q(x):=B(x,x)$.
		
		\textbf{Dfn:} For any quadratic form $Q$ the associated symmetric bi-linear form is given by the \textbf{polarization identity}:
		$
		B(x,y)=\frac{1}{2}(Q(x+y)-Q(x)-Q(y)).
		$
		
		\textbf{Dfn:} A norm on $\mathbb{R}^n$ is a map $Q:\mathbb{R}^n\to\mathbb{R}$ which obeys $Q(x)\geq 0$ for all $x\in\mathbb{R}^n$ with $Q(x)=0$ iff $x=0$.
		
		\textbf{Lem:} (13.2). A symmetric bi-linear form is positive-definite iff its associated quadratic form defines a norm on $\mathbb{R}^n$.
		
		\textbf{Dfn:} (2$^{nd}$ Variation). The 2$^{nd}$ variation of an action functional $S[x]$ is the quadratic form
		$
		Q^S_x(\varepsilon):=\left.\frac{d^2}{ds^2}S[x+s\varepsilon]\right|_{s=0},
		$
		acting on the space of variations $\varepsilon:[0,1]\to\mathbb{R}^n$ with appropriate boundary conditions (e.g.\ endpoint fixed variations).
		
		Making use of the action in terms of the Lagrangian and restricting our attention to endpoint fixed variations,
		$
		Q^S_x[\varepsilon]=\int_0^1\sum_{i,j=1}^n\left(\frac{\partial L}{\partial x_i x_j}\varepsilon_i\varepsilon_j+2\frac{\partial L}{\partial\dot{x}_i x_j}\dot\varepsilon_i\varepsilon_j+\frac{\partial L}{\partial\dot{x}_i\dot{x}_j}\dot\varepsilon_i\dot\varepsilon_j\right)dt.
		$
		For $n=1$;
		$
		Q^S_x[\varepsilon]=\int_0^1\left(\frac{\partial^2 L}{\partial x^2}\varepsilon^2+2\frac{\partial L}{\partial x,\dot{x}}\varepsilon\dot\varepsilon+\frac{\partial^2 L}{\partial\dot{x}^2}\dot\varepsilon^2\right)dt.
		$
		Then the polarization identity (by IBP with $\varepsilon(0)=\varepsilon(1)=0$ and $\eta(0)=\eta(1)=0$ gives
		$
		H^S_x(\varepsilon,\eta)=\int_0^1\eta\left(\frac{\partial^2 L}{\partial x^2}-\frac{d}{dt}\left(\frac{\partial L}{\partial x,\dot{x}}\right)-\frac{d}{dt}\left(\frac{\partial L}{\partial\dot{x}}\right)\frac{d}{dt}\right)\varepsilon dt.
		$
		Which can be written compactly as $H^S_x(\varepsilon,\eta)=\int_0^1\eta J(\varepsilon)dt$, where
		$
		J=\frac{\partial^2 L}{\partial x^2}-\frac{d}{dt}\left(\frac{\partial L}{\partial x,\dot{x}}\right)-\frac{d}{dt}\left(\frac{\partial L}{\partial\dot{x}}\right)\frac{d}{dt},
		$
		is a 2$^{nd}$-order differential operator in $t$. This is generalised into $n>1$ dims by
		$
		H^S_x(\varepsilon,\eta)=\int_0^1\left(\sum_{i,j=1}^n\eta_i J_{ij}(\varepsilon_j)\right)dt=\int_0^1\eta\cdot J(\varepsilon)dt,
		$
		where
		$
		J_{ij}:=\frac{\partial L}{\partial x_i x_j}-\frac{1}{2}\frac{d}{dt}\left(\frac{\partial L}{\partial\dot{x}_i x_j}+\frac{\partial L}{\partial x_i\dot{x}_j}\right)-\frac{d}{dt}\left(\frac{\partial L}{\partial\dot{x}_i\dot{x}_j}\frac{d}{dt}\right)
		$
		define the components of the $n\times n$ matrix of 2$^{nd}$-order differential operators $J$. This matrix is called the \textbf{Jacobi operator} of the Lagrangian.
		
		A soln $x$ to the Euler-Lagrange eqns is a (local) min of the action $S[x]$ iff $H^S_x$ is positive-definite on the space of endpoint fixed variations, or, equivalently if $Q^S_x(\varepsilon)>0$ for all non-trivial endpoint fixed variations $\varepsilon$. Likewise, $x$ is a (local) max iff $H^S_x$ is negative-definite or if $Q^S_x(\varepsilon)<0$ for all non-trivial endpoint fixed variations.
		
		\hrule
		\vspace*{-3pt}
		\subsubsection*{\tiny{14. Principle of Least Action}}
		\vspace*{-5pt}
		
		\textbf{Princ:} (Hamilton's Principle of Least Action). Recall that Newt's eqns for a conservative mechanical system of the general form is
		$
		M\ddot{x}=-\nabla U(x)
		$
		where the dim of configuration space is $N$, $M$ is the symmetric $N\times N$ mass matrix of the system and $U(x)$ is the potential. Hamilton's principle of least action states that these eqns always arise as the extrema of an action functional whose Lagrangian takes a particular form.
		
		\textbf{Thm:} (14.1). Motion for a mechanical system coincides with the extremals of the action functional with Lagrangian
		$
		L(x,\dot{x};t)=T-U
		$
		where $T$ is kinetic energy and $U$ is potential.
		
		\textbf{Lem:} (14.1). For the Lagrangian $L=T-U$ of a conservative mechanical system, the energy is $E=T+U$ agreeing with the physical definition of total energy of the mechanical system.
		
		\textbf{Ex:} (Free Particle). Write the Lagrangian for a free particle of mass $m$ and show that the solns of the E-L eqns minimize the action. For a free particle the only energy is kinetic energy so $U=0$ and
		$
		T=\frac{m}{2}|\dot{x}|^2, x:[0,1]\to\mathbb{R}^3
		$
		with $L=T=\frac{m}{2}|\dot{x}|^2$. The E-L eqns are simply
		$
		\frac{d}{dt}=m\ddot{x}=0
		$
		whose solns are straight lines. Note that since $\frac{\partial L}{\partial x}=0$ the momentum $p=\frac{\partial L}{\partial\dot{x}}=m\dot{x}$ is cnsrvd along with energy. Calculating the 2$^{nd}$ variation of the action:
		$
		Q^S_x(\varepsilon)=m\int_0^1|\dot\varepsilon|^2dt>0\text{ for all nontrivial }\varepsilon,
		$
		so the solns of the E-L eqns are the minima of the action.
		
		\textbf{Ex:} (Pendulum with a moving pivot). Consider a planar pendulum whose pivot is a body of mass $M$ which can move freely in the horizontal direction. Write the Lagrangian and E-L eqns for this system.
		
		Place the origin on the horizontal line along which the mass $M$ can move, and let $x$ denote the horizontal displacement of the mass $M$ from this point. Let
		$
		x_2=\begin{psmallmatrix*}x+L\sin\theta\\-L\cos\theta\end{psmallmatrix*}
		$
		denote the position of the mass $m$ at the end of the pendulum, relative to origin. Then
		$
		|\dot{x}_2|^2=\dot{x}^2+2L\cos\theta\dot{x}\dot\theta+L^2\dot\theta^2,
		$
		so the kinetic energy is given by
		$
		T=\frac{M}{2}\dot{x}^2+\frac{m}{2}|\dot{x}_2|^2=\frac{(M+m)}{2}\dot{x}^2+\frac{m}{2}L^2\dot\theta^2+mL\cos\theta\dot{x}\dot\theta.
		$
		This gives the mass matrix as
		$
		\begin{psmallmatrix*}M+m & mL\cos\theta\\ mL\cos\theta & mL^2\end{psmallmatrix*}
		$
		which is not diagonal! The only potential from the system is gravitational from the pendulum which is $-mgL(1-\cos\theta)$. The Lagrangian is thus
		$
		\mathcal{L}=\frac{M}{2}\dot{x}^2+\frac{m}{2}|\dot{x}_2|^2=\frac{(M+m)}{2}\dot{x}^2+\frac{m}{2}L^2\dot\theta^2+mL\cos\theta\dot{x}\dot\theta
		-mgL(1-\cos\theta).
		$
		Note that as $\frac{\partial\mathcal{L}}{\partial x}=0$ the momentum
		$
		p=(M+m)\dot{x}+mL\cos\theta\dot\theta
		$
		is cnsrvd, in addition to the total energy. The E-L eqns for the system are
		$
		\frac{d}{dt}\left(\dot{x}+\frac{m}{M}\frac{L\cos\theta\dot\theta}{1+\frac{m}{M}}\right)=0,
		$
		$
		-\sin\theta\left(\frac{g}{L}+\frac{\dot{x}\dot\theta}{L}\right)=\frac{d}{dt}\left(\dot\theta+\frac{\cos\theta\dot{x}}{L}\right).
		$
		
		\hrule
		\vspace*{-3pt}
		\subsubsection*{\tiny{15. Noether's Thm}}
		\vspace*{-5pt}
		
		\textbf{Diffeomorphisms and symmetries:}
		
		Let $S[x]$ be a functional, assume $x\in C^1([0,1],\mathbb{R}^n)$ and $x(0)=P$ and $x(1)=Q$ for $P,Q\in\mathbb{R}^n$. Consider a map $\varphi:\mathbb{R}^n\to\mathbb{R}^n$ which is twice cont.ly diffable. It follows that
		$
		y:=\varphi\circ x, y(t)=\varphi(x(t))
		$
		is also a member of $C^1([0,1],\mathbb{R}^n)$ with endpoints $y(0)=\varphi(P)$ and $y(1)=\varphi(Q)$ in $\mathbb{R}^n$.
		
		\textbf{Dfn:} ($C^2$ Diffeomorphism). A map $\varphi:\mathbb{R}^n\to\mathbb{R}^n$ is called a $C^2$ diffeomorphism if $\varphi$ is invertible and both $\varphi$ and $\varphi^{-1}$ are in $C^2(\mathbb{R}^n,\mathbb{R}^n)$.
		
		Now let $L:\mathbb{R}^{2n+1}\to\mathbb{R}$ be any smooth Lagrangian and define $K:\mathbb{R}^{2n+1}\to\mathbb{R}$ by
		$
		K(x,\dot{x};t):=L(\varphi(x),D\varphi(x)\dot{x};t)\text{where}
		$
		$
		D\varphi(x):=\begin{psmallmatrix*}\frac{\partial\varphi}{\partial x_1}&\cdots&\frac{\partial\varphi}{\partial x_n}\end{psmallmatrix*}
		$
		is the Jacobian of partial derivatives of $\varphi$.
		
		\textbf{Lem:} (15.1). Let $S[y]=\int_0^1 L(y,\dot{y};t)dt$, $I[x]=\int_0^1 K(x,\dot{x};t)dt$ then $x$ extremises $I$ iff $y=\varphi\circ x$ extremises $S$.
		
		\textbf{Dfn:} (Symmetry and Invariance). If a $C^2$ diff.ism $\varphi:\mathbb{R}^n\to\mathbb{R}^n$ has $L(x,\dot{x};t)=L(y,\dot{y};t)\text{for }y=\varphi\circ x.$, it is a \textbf{symmetry} of $L$/$L$ is \textbf{invariant} under $\varphi$.
		
		\textbf{Cor:} (15.1). If $\varphi$ is a symm of the Lagrangian, then $\varphi$ maps extrema of the action to extrema of the action.
		
		Consider a 1 param family of maps $\{\varphi_s:\mathbb{R}^n\to\mathbb{R}^n|s\in\mathbb{R}\}$ subject to the conditions:
		\vspace*{-10pt}
		\begin{multicols}{2}
		\begin{enumerate}[leftmargin=*,topsep=0pt,itemsep=0pt]
			\item $\forall s\in\mathbb{R}$, $\varphi_s$ is a $C^2$ func on $\mathbb{R}^n$,
			\item $\varphi_s(x)$ is diffable in $s\in\mathbb{R}$,
			\item $\varphi_0(x)=x$,
			\item $\forall s,t\in\mathbb{R}$, $\varphi_s\circ\varphi_t=\varphi_{s+t}$.
		\end{enumerate}
		\end{multicols}
		\vspace*{-10pt}
		These conditions ensure that $\varphi_s$ is invertible with $\varphi_s^{-1}=\varphi_{-s}$. Thus the 1 param family is actually a 1 param group of $C^2$ diff.isms on $\mathbb{R}^n$.
		
		\textbf{Thm:} (Noether's Thm). Let $S[x]=\int_0^1 L(x,\dot{x};t)dt$ be an action for curves $x:[0,1]\to\mathbb{R}^n$ and let $L$ be invariant under the one param group of $C^2$ diff.isms $\{\varphi_s\}$. Then the Noether charge
		$
		N(x,\dot{x};t):=\frac{\partial L}{\partial\dot{x}}\cdot\frac{\partial\varphi_s}{\partial s}\bigg|_{s=0}
		$
		is a cnsrvd qtty along solns of the E-L eqns.
		
		\textbf{Cor:} (15.2). Suppose a 1 param group of $C^2$ diff.isms $\{\varphi_s\}$ transform the Lagrangian s.t.
		$
		\frac{\partial L}{\partial s}(y,\dot{y};t)=\frac{\partial K_s}{\partial t}(x,\dot{x};t)\text{for }y=\varphi_s\circ x
		$
		and some $K_s(x,\dot{x};t)$. Then the Noether charge
		$
		N(x,\dot{x};t)=\frac{\partial L}{\partial\dot{x}}\cdot\frac{\partial\varphi_s}{\partial s}\bigg|_{s=0}-K_0(x,\dot{x};t)
		$
		is a cnsrvd qtty along solns of the E-L eqns.
		
		\textbf{Thm:} (Noether's Thm -- Infinitesimal). Let $\{\varphi_s\}$ be a one param family of $C^2$ diff.isms and let
		$
		\zeta(X):=\frac{\partial\varphi_s}{\partial s}\bigg|_{s=0}.
		$
		If the Lagrangian $L$ obeys
		$
		\frac{\partial L}{\partial x}\cdot\zeta+\frac{\partial L}{\partial\dot{x}}\cdot\dot\zeta=\frac{dK}{dt},
		$
		then the Noether charge
		$
		N(x,\dot{x};t)=\frac{\partial L}{\partial\dot{x}}\cdot\zeta-K
		$
		is a cnsrvd qtty along solns of E-L eqns.
		
		\hrule
		\vspace*{-3pt}
		\subsubsection*{\tiny{16. Generalising Noether's Thm + Exs}}
		\vspace*{-5pt}
		
		\textbf{Dfn:} (Non-autonomous system). A non-autonomous Lagrangian system is one for which the kinetic and potential energies have the functional dependencies $T=T(x,\dot{x};t)$ and $U=U(x;t)$ respectively. Equivalently, the Lagrangian $L=T-U$ has explicit time dependence.
		
		Let $\{\psi_s:\mathbb{R}^{n+1}\to\mathbb{R}^{n+1}\}$ be a one param family of $C^2$ diff.isms for $s\in\mathbb{R}$ which act as
		$
		\psi_s(x,t)=(y,\bar{t})\text{subject to}
		$
		$
		y=x+s\zeta(x,t)+O(s^2),\bar{t}=t+s\tau(x,t)+O(s^2)\text{for}
		$
		$
		\zeta:=\frac{\partial y}{\partial s}\bigg|_{s=0},\tau:=\frac{\partial\bar{t}}{\partial s}\bigg|_{s=0}.
		$
		This expansion allows us to is off a variety of useful partial derivatives:
		$
		\frac{\partial y_j}{\partial x_k}\bigg|_{s=0}=\delta_{jk}, \frac{\partial y_j}{\partial t}\bigg|_{s=0}=0, \frac{\partial\bar{t}}{\partial x_k}\bigg|_{s=0}=0, \frac{\partial\bar{t}}{\partial t}\bigg|_{s=0}=1
		$
		$
		\frac{\partial y_j}{\partial s,x_k}\bigg|_{s=0}=\frac{\partial\zeta_j}{\partial x_k}, \frac{\partial y_j}{\partial s,t}\bigg|_{s=0}=\frac{\partial\zeta_j}{\partial t}, \frac{\partial\bar{t}}{\partial s,x_k}\bigg|_{s=0}=\frac{\partial\tau}{\partial x_k}, \frac{\partial\bar{t}}{\partial s,t}\bigg|_{s=0}=\frac{\partial\tau}{\partial t}.
		$
		This allows us to define the notion of a symm for non-autonomous systems.
		
		\textbf{Dfn:} A 1-param family of $C^2$ diff.isms $\{\psi_s(y,\bar{t})\}$ is a symm of the Lagrangian if
		$
		L\!\left(y,\frac{dy}{d\bar{t}};\bar{t}\right)\frac{d\bar{t}}{dt}=L(x,\dot{x};t).
		$
		
		\textbf{Thm:} (Noether's Thm for non-autonomous systems). Let the 1-param family of $C^2$ diff.isms $\{\psi_s=(y,\bar{t})\}$ be an infinitesimal symm of the Lagrangian, in the sense that
		$
		\frac{\partial L}{\partial x}\cdot\zeta+\frac{\partial L}{\partial\dot{x}}\cdot\frac{d\zeta}{dt}=\frac{d}{dt}\!\left(\frac{\partial L}{\partial\dot{x}}\cdot\zeta\right)+\left(\frac{\partial L}{\partial x}-\frac{d}{dt}\frac{\partial L}{\partial\dot{x}}\right)\cdot\zeta
		$
		holds. Then the Noether charge
		$
		N(x,\dot{x};t):=L\tau+\frac{\partial L}{\partial\dot{x}}\cdot(\zeta-\dot{x}\tau)
		$
		is a cnsrvd qtty along solns of the E-L eqns.
		
		\textbf{Lem:} (16.1). A 1-param family of $C^2$ diff.isms is an infinitesimal symm of the non autonomous Lagrangian if
		$
		\frac{\partial L}{\partial t}\tau+\frac{\partial L}{\partial x}\cdot\zeta+\frac{\partial L}{\partial\dot{x}}\cdot\left(\frac{\partial\zeta}{\partial t}-\dot{x}\frac{\partial\tau}{\partial t}\right)+L\frac{\partial\tau}{\partial t}=0.
		$
		
		\textbf{Cor:} (16.1). If a one-param family of $C^2$ diff.isms $\{\psi_s=(y,\bar{t})\}$ act infinitesimally on the Lagrangian as
		$
		\frac{\partial L}{\partial t}\tau+\frac{\partial L}{\partial x}\cdot\zeta+\frac{\partial L}{\partial\dot{x}}\cdot\left(\frac{d\zeta}{dt}-\dot{x}\frac{d\tau}{dt}\right)+L\frac{d\tau}{dt}=\frac{dK}{dt}(x,\dot{x};t),
		$
		for some $K(x,\dot{x};t)$, then the Noether charge
		$
		N(x,\dot{x};t):=L\tau+\frac{\partial L}{\partial\dot{x}}\cdot(\zeta-\dot{x}\tau)-K
		$
		is a cnsrvd qtty along solns of the E-L eqns.
		
		\textbf{Ex:} (Linear Momentum). Consider the Lagrangian $L=\frac{1}{2}\dot{x}\cdot M\dot{x}-U(x)$, where $x:[0,1]\to\mathbb{R}^n$ and $M$ is the $n\times n$, real valued, symmetric mass matrix. Suppose that $U(x)$ does not depend on $x_i$ for some $i=1,\ldots,n$. Show that $\varphi_s(x)$, defined by
		$
		\varphi_s(x)= x_j+s \text{ if }j=i$ and $=x_j\text{ otherwise }
		$
		is a symm of the Lagrangian and deduce the associated cnsrvd qtty. As the Lagrangian is indep. of $x_i$ and $\dot\varphi=\dot{x}$, it follows immediately that $\varphi_s$ is a symm of the Lagrangian. It is also clearly a $C^2$ diff.ism and so by Noether's thm the Noether charge is
		$
		N=\frac{\partial L}{\partial\dot{x}}\cdot\frac{\partial\varphi_s}{\partial s}\bigg|_{s=0}=\sum_{j,k=1}^n M_{jk}\dot{x}_k\delta_{ij}=M_{jk}\dot{x}_k,
		$
		which is the generalised momentum component $p_i$.
		
		\hrule
		\vspace*{-3pt}
		\subsubsection*{\tiny{17. Hamiltonian Formalism}}
		\vspace*{-5pt}
		
		\textbf{Dfn:} (Regular Lagrangian). A Lagrangian $L(x,\dot{x};t)$ is \textbf{regular}, or non-degen., if the $n\times n$ matrix with entries
		$
		\frac{\partial L}{\partial\dot{x}_i\partial\dot{x}_j}, i,j=1,\ldots,n,
		$
		is invertible for all values of $(x,\dot{x},t)\in\mathbb{R}^{2n+1}$.
		
		\textbf{Dfn:} (Hamiltonian). Let $L(x,\dot{x};t)$ be a regular Lagrangian. The func
		$
		H(x,p;t):=p\cdot\dot{x}-L(x,\dot{x};t),\text{for }\dot{x}=\dot{x}(x,p;t),
		$
		is called the \textbf{Hamiltonian}.
		
		\textbf{Thm:} (17.1, Hamilton's Eqns). The Euler-Lagrange eqns of a regular Lagrangian system are equivalent to the system of $2n$ 1$^{st}$-order differential eqns
		$
		\frac{dx}{dt}=\frac{\partial H}{\partial p}, \frac{dp}{dt}=-\frac{\partial H}{\partial x},
		$
		where $H=H(x,p;t)$ is the Hamiltonian. These eqns are called \textbf{Hamilton's Eqns} or the canonical form for the Euler-Lagrange eqns.
		
		\textbf{Dfn:} (Phase Space). $\mathbb{R}^{2n}$ with coordinates $(x,p)$ is the \textbf{phase space} of the mechanical system.
		
		Upon changing variables from the state space $(x,\dot{x})$ to the phase space $(x,p)$ the total energy of a mechanical system is seen to be equal to the Hamiltonian:
		$
		E=\dot{x}\cdot\frac{\partial L}{\partial\dot{x}}-L=\dot{x}(x,p;t)\cdot p-L(x,\dot{x};t)=H.
		$
		
		\textbf{Dfn:} (Poisson bracket). Let $f,g:\mathbb{R}^{2n}\to\mathbb{R}$ be any two diffable funcs on the phase space. The \textbf{Poisson bracket} of $f$ and $g$ is defined by
		$
		\{f,g\}:=\frac{\partial f}{\partial x}\cdot\frac{\partial g}{\partial p}-\frac{\partial f}{\partial p}\cdot\frac{\partial g}{\partial x}.
		$
		
		\textbf{Lem:} (17.2). The time-evolution of any qtty $\Phi:\mathbb{R}^{2n}\to\mathbb{R}$ (that does not explicitly depend on time ($\frac{\partial\Phi}{\partial t}=0$)) on the phase space along solns to the eqns of motion is governed by the Poisson bracket with the Hamiltonian
		$
		\frac{d\Phi}{dt}=\{\Phi,H\}.
		$
		Furthermore, $\Phi$ is a cnsrvd qtty iff it `Poisson commutes' with the Hamiltonian, $\{\Phi,H\}=0$.
		
		\textbf{Prop:} (17.1). The Poisson bracket has the following properties:
		\vspace*{-10pt}
		\begin{multicols*}{2}
		\begin{itemize}[leftmargin=*,topsep=0pt,itemsep=0pt]
			\item $\{\cdot,\cdot\}$ is bi-linear in its arguments.
			\item $\{f,g\}=-\{g,f\}$.
			\item $\{f,gh\}=\{f,g\}h+g\{f,h\}$.
		\end{itemize}
	\end{multicols*}
	\vspace*{-10pt}
	\begin{itemize}
		\item Jacobi identity: $\{f,\{g,h\}\}=\{\{f,g\},h\}+\{g,\{f,h\}\}$
	\end{itemize}
		for all twice-diffable funcs $f,g,h:\mathbb{R}^{2n}\to\mathbb{R}$ on the phase space.
		
		\textbf{Lem:} (17.3). If $\Phi_1(x,p)$ and $\Phi_2(x,p)$ are cnsrvd quantities, then so is $\{\Phi_1,\Phi_2\}$.
		
		\textbf{Thm:} (17.2). Every cnsrvd qtty in a mechanical system corresponds to a one-param family of symmetries of that system.
		
		This is not the exactly converse of Noether's thm(s), which require a one-param family of symmetries which are $C^2$ diff.isms. Whereas the proof above only implicitly assumes that the one-param family is cont.ly diffable.
		
		\hrule
		\vspace*{-3pt}
		\subsubsection*{\tiny{18. Canonical Transformations and Integrability}}
		\vspace*{-5pt}
		
		\textbf{Dfn:} (Canonical transform). Let $(x,p)$ be Darboux coordinates on phase space. A change of variables $(x,p)\mapsto(X,P)$ is called a \textbf{canonical transform} if the Poisson brackets of $(X,P)$, calculated as funcs of $(x,p)$, obey
		$
		\{X_i,X_j\}=0=\{P_i,P_j\},\{X_i,P_j\}=\delta_{ij}.
		$
		
		\textbf{Thm:} Let $(x,p)$ be Darboux coordinates on the phase space of a system subject to Hamilton's eqns with Hamiltonian $H(x,p,t)$. Let $x,p\mapsto(X,P)$ be a canonical transform and let $\tilde{H}$ be the Hamiltonian in the new variables. Then the eqns of motion in the new variables are
		$
		\frac{dX}{dt}=\frac{\partial\tilde{H}}{\partial P},\frac{dP}{dt}=-\frac{\partial\tilde{H}}{\partial X}.
		$
		
		$\Phi_1$ and $\Phi_2$ are \textbf{independent} if the vectors $\nabla\Phi_1,\nabla\Phi_2\in\mathbb{R}^{2n}$ are lin. ind. for all values of $(x,p)$.
		
		\textbf{Dfn:} (Integrable system). A mechanical system whose phase space is $2n$-dim.al is (Hamiltonian) \textbf{integrable} if it has $n$ indep. cnsrvd quantities, $\Phi_1,\Phi_2,\ldots,\Phi_n$, with $\Phi_1=H$ the Hamiltonian of the system, all of which Poisson commute:
		$
		\{\Phi_i,\Phi_j\}=0,\text{for all }i,j,\ldots,n.
		$
		
		\textbf{Lem:} Every 1D mechanical system is integrable.
		
		\hrule
		\vspace*{-3pt}
		\subsubsection*{\tiny{19. Isoperimetric Problems}}
		\vspace*{-5pt}
		
		\textbf{Lem:} Let $x\in U\subseteq\mathbb{R}^n$ be an extremum of $f$ subject to $g=0$. Then if $\nabla g(x)\neq 0$, there exists a Lagrange multiplier $\lambda\in\mathbb{R}$ s.t. $(x,\lambda)$ is a CP of the func $F:U\times\mathbb{R}\to\mathbb{R}$ given by $F:=f-\lambda g$.
		
		\textbf{Thm:} Suppose that $x(t)\in C_{P,Q}$ be an extremum of $S$ subject to $I[x]=0$. Then if $x(t)$ is not an extremum of $I$, there exists a Lagrange multiplier $\lambda\in\mathbb{R}$ s.t. $x(t)$ is an extremum of the functional $S[x]-\lambda I[x]$.
		
		\textbf{Dfn:} (Method of Lagrange multipliers).
		\begin{enumerate}[leftmargin=*,topsep=0pt,itemsep=0pt]
			\item Ensure that $I[x]$ has no extrema obeying $I[x]=0$.
			\item Solve the Euler-Lagrange eqns of $L-\lambda K$.
			\item Fix constants of integration using the boundary conditions $x(0)=P$, $x(1)=Q$.
			\item Fix the value of $\lambda$ from $I[x]=0$.
		\end{enumerate}
		This only gives a necessary condition for $x$ to be a soln of the isoperimetric problem. Normally ok to manually check the resulting $x$ is a soln. Just add a new Lagrange multiplier for each new constraint in the problem.
		
		\hrule
		\vspace*{-3pt}
		\subsubsection*{\tiny{20. Holonomic/Nonholonomic Constraints}}
		\vspace*{-5pt}
		
		\textbf{Dfn:} (Constraints). To determine extrema of an action $S[x]$ for $x\in C_{P,Q}$ subject to a constraint $G(x,\dot{x},t)=0$, note there are 4 different types:
		\begin{enumerate}[leftmargin=*,topsep=0pt,itemsep=0pt]
			\item Scleronomic; if $G$ has no explicit time dependence.
			\item Rheonomic if $G$ depends explicitly on time.
			\item Holonomic if $G$ does not depend explicitly on $\dot{x}$.
			\item Non-holonomic if $G$ has explicit $\dot{x}$ dependence.
		\end{enumerate}
		
		Let $S[x]$ be subject to a holonomic constraint $G(x,\dot{x},t)=g(x,t)=0$ for $g:\mathbb{R}^n\times[0,1]\to\mathbb{R}$ smooth and $x\in C_{P,Q}$. The boundary conditions on $x$ imply $g(P,0)=0=g(Q,0)$, we assume $\nabla g(x,t)\neq 0\forall t\in[0,1]$.
		
		
		\textbf{Thm:} (20.1). Let $x\in C_{P,Q}$ be a smooth extremum of $S[x]$ subject to the holonomic constraint $g(x,t)=0$, with $\nabla g(x,t)\neq 0$ for all $t\in[0,1]$. Then there exists some $\lambda:[0,1]\to\mathbb{R}$ s.t. $x$ obeys the E-L eqns of $L-\lambda g$.
		
		This means the method of Lagrange multipliers now extends to the setting of holonomic constraints:
		\begin{enumerate}[leftmargin=*,topsep=0pt,itemsep=0pt]
			\item Establish that $\nabla g(x,t)\neq 0$.
			\item Solve the E-L eqns of $L-\lambda g$.
			\item Fix constants of integration with boundary conditions.
			\item Determine $\lambda(t)$ by imposing $g(x(t),t)=0$.
		\end{enumerate}
		
		\textbf{Dfn:} A constraint (non-holo) $G(x,\dot{x},t)=0$ is called a \textbf{Pfaffian constraint} if it is at most linear in $\dot{x}$.
		
		\textbf{Lem:} (20.1). Let $x\in C_{P,Q}$ be a smooth extremum of $S[x]$ subject to a Pfaffian constraint $G(x,\dot{x},t)=0$ with $\nabla G(x,\dot{x},t)\neq 0$ for all $t\in[0,1]$. Then there exists some $\lambda:[0,1]\to\mathbb{R}$ s.t. $x$ obeys the E-L eqns $L-\lambda G$.
		
	\end{multicols*}
\end{document}
